\documentclass[12pt,a4paper]{article}

\usepackage[utf8]{inputenc}
\usepackage[T2A]{fontenc}
\usepackage[russian]{babel}
\usepackage{amsmath, amsfonts, amssymb, amsthm}
\usepackage{graphicx}
\usepackage{tikz}
\usetikzlibrary{shapes.geometric, arrows, positioning}
\usepackage{algorithm}
\usepackage{algpseudocode}
\usepackage{geometry}
\geometry{top=2cm, bottom=2cm, left=2.5cm, right=1.5cm}

\newtheorem{theorem}{Теорема}
\renewcommand{\proofname}{Доказательство}
\newtheorem{definition}{Определение}

\floatname{algorithm}{Алгоритм}

\title{Билет №82 \\ \small Экспертные системы (ЭС). Архитектура ЭС. Механизмы вывода, подсистемы объяснения, общения, приобретения знаний ЭС. Жизненный цикл экспертной системы. (ИТМО)}
\author{Сериков Владислав Александрович}
\date{16 мая 2026}

\begin{document}

\maketitle
\tableofcontents
\newpage

\section{Введение}

\begin{definition}
\textbf{Экспертная система (ЭС)} --- это вычислительная система, в которую включены знания специалистов о некоторой узкой предметной области и которая в пределах этой области способна принимать экспертные решения (диагностика, проектирование, планирование).
\end{definition}

Главное отличие ЭС от традиционных программ заключается в четком разделении \textit{базы знаний} (декларативных или процедурных знаний о предметной области) и \textit{механизма вывода} (алгоритма применения этих знаний).

\section{Архитектура экспертных систем}

Классическая экспертная система состоит из следующих взаимосвязанных компонентов (рис. \ref{fig:architecture}):

\begin{figure}[htbp]
\centering
\begin{tikzpicture}[
    node distance=2cm, auto, thick,
    box/.style={draw, rectangle, minimum height=1cm, minimum width=2.5cm, align=center},
    arrow/.style={->, >=stealth, thick}
]

\node[box] (User) {Пользователь};
\node[box, right=1.5cm of User] (Interface) {Подсистема \\ общения \\ (Интерфейс)};
\node[box, right=1.5cm of Interface] (Inference) {Механизм \\ логического \\ вывода};
\node[box, above=1cm of Inference] (Explanation) {Подсистема \\ объяснения};
\node[box, below=1cm of Inference] (WorkingMem) {Рабочая \\ память (БД фактов)};
\node[box, right=1.5cm of Inference] (KB) {База \\ знаний};
\node[box, below=1cm of KB] (Acquisition) {Подсистема \\ приобретения \\ знаний};
\node[box, right=1.5cm of Acquisition] (Expert) {Инженер по знаниям \\ / Эксперт};

\draw[<->, >=stealth, thick] (User) -- (Interface);
\draw[<->, >=stealth, thick] (Interface) -- (Inference);
\draw[<->, >=stealth, thick] (Interface) |- (Explanation);
\draw[<->, >=stealth, thick] (Inference) -- (WorkingMem);
\draw[<->, >=stealth, thick] (Inference) -- (KB);
\draw[<->, >=stealth, thick] (Inference) -- (Explanation);
\draw[<->, >=stealth, thick] (Acquisition) -- (KB);
\draw[<->, >=stealth, thick] (Expert) -- (Acquisition);

\end{tikzpicture}
\caption{Базовая архитектура экспертной системы}
\label{fig:architecture}
\end{figure}

\begin{enumerate}
    \item \textbf{База знаний (БЗ)} --- хранилище долгосрочных знаний о предметной области (правила, фреймы, семантические сети, онтологии).
    \item \textbf{Рабочая память (РП) / База данных} --- хранилище краткосрочных данных. Содержит текущие факты, полученные от пользователя или выведенные системой в ходе сеанса.
    \item \textbf{Механизм логического вывода (МЛВ)} --- программа, моделирующая ход рассуждений эксперта на основании знаний из БЗ и фактов из РП.
    \item \textbf{Подсистема объяснения} --- модуль, аргументирующий выводы системы (ответы на вопросы «Почему?» и «Как?»).
    \item \textbf{Подсистема общения (интерфейс)} --- модуль взаимодействия системы с конечным пользователем.
    \item \textbf{Подсистема приобретения знаний} --- модуль для автоматизации наполнения БЗ (через диалог с инженером по знаниям, экспертом или методами машинного обучения).
\end{enumerate}

\section{Механизмы логического вывода}

Фундаментом МЛВ в продукционных (основанных на правилах) ЭС является формальная логика. Вывод опирается на теоремы математической логики, гарантирующие корректность получаемых результатов.

\begin{theorem}[Корректность правила резолюции]
Пусть $C_1 = A \vee L$ и $C_2 = B \vee \neg L$ --- два дизъюнкта, где $L$ --- литерал. Если дизъюнкт $C = A \vee B$ (резольвента) получен из $C_1$ и $C_2$ по правилу резолюции, то он логически следует из них: $C_1, C_2 \models C$.
\end{theorem}

\begin{proof}
Предположим, что существует некоторая интерпретация $I$ (набор истинностных значений), при которой $C_1$ и $C_2$ истинны. Возможны два случая для литерала $L$:
\begin{enumerate}
    \item В интерпретации $I$ литерал $L$ истинен ($L = True$). Тогда $\neg L$ ложен ($False$). Так как $C_2 = B \vee \neg L$ истинно по предположению, единственным способом сделать дизъюнкцию $C_2$ истинной является истинность $B$. Если $B$ истинно, то и дизъюнкция $C = A \vee B$ истинна в интерпретации $I$.
    \item В интерпретации $I$ литерал $L$ ложен ($L = False$). Тогда дизъюнкция $C_1 = A \vee L$ может быть истинной только при истинности $A$. Если $A$ истинно, то и дизъюнкция $C = A \vee B$ истинна в $I$.
\end{enumerate}
В обоих возможных случаях при истинности посылок истинно и следствие. Таким образом, $C$ логически следует из $C_1$ и $C_2$. Теорема доказана.
\end{proof}

На практике МЛВ чаще всего использует два основных алгоритма (направления вывода): \textbf{прямой} и \textbf{обратный}.

\subsection{Прямой вывод (Forward Chaining)}

Прямой вывод реализует стратегию «от фактов к цели». Система сопоставляет известные факты из Рабочей памяти с условиями (антецедентами) правил из БЗ.

\begin{algorithm}[H]
\caption{Алгоритм прямого логического вывода}
\begin{algorithmic}[1]
\Require Рабочая память $WM$ (начальные факты), База правил $R$
\State $GoalReached \gets \text{false}$
\While{не $GoalReached$ \textbf{and} существуют применимые правила}
    \State $ConflictSet \gets \emptyset$
    \For{каждого правила $r \in R$}
        \If{все условия $r$ присутствуют в $WM$}
            \State $ConflictSet \gets ConflictSet \cup \{r\}$
        \EndIf
    \EndFor
    \If{$ConflictSet \neq \emptyset$}
        \State $r_{selected} \gets \text{ResolveConflict}(ConflictSet)$ \Comment{Разрешение конфликта}
        \State $WM \gets WM \cup \{\text{заключение } r_{selected}\}$
        \State $R \gets R \setminus \{r_{selected}\}$
        \If{заключение $r_{selected}$ является целевым фактом}
            \State $GoalReached \gets \text{true}$
        \EndIf
    \Else
        \State \textbf{break} \Comment{Нет применимых правил}
    \EndIf
\EndWhile
\Return $WM$
\end{algorithmic}
\end{algorithm}

\textbf{Пример прямого вывода:} \\
\textit{Правила:} \\
$R_1$: ЕСЛИ "двигатель не заводится" И "стартер крутит" ТО "проверить систему зажигания". \\
$R_2$: ЕСЛИ "проверить систему зажигания" И "искры нет" ТО "неисправна катушка зажигания". \\
\textit{Факты в $WM$:} "двигатель не заводится", "стартер крутит", "искры нет". \\
\textit{Шаг 1:} Условия $R_1$ выполнены. Добавляем в $WM$ факт "проверить систему зажигания". \\
\textit{Шаг 2:} Условия $R_2$ теперь выполнены. Добавляем "неисправна катушка зажигания". Цель достигнута.

\subsection{Обратный вывод (Backward Chaining)}

Обратный вывод реализует стратегию «от цели к фактам». Система выдвигает гипотезу (цель) и пытается найти правила, подтверждающие её, рекурсивно доказывая их посылки.

\begin{algorithm}[H]
\caption{Алгоритм обратного логического вывода}
\begin{algorithmic}[1]
\Require Цель $G$, Рабочая память $WM$, База правил $R$
\Function{Prove}{$G$}
    \If{$G \in WM$}
        \Return \text{true}
    \EndIf
    \State $RelevantRules \gets \{r \in R \mid \text{заключение } r = G\}$
    \If{$RelevantRules = \emptyset$}
        \State \textbf{AskUser}(«Истинно ли $G$?») \Comment{Запрос факта у пользователя}
        \Return ответ пользователя
    \EndIf
    \For{каждого правила $r \in RelevantRules$}
        \State $RuleProved \gets \text{true}$
        \For{каждого подусловия $P_i$ в условии правила $r$}
            \If{\textbf{not} \Call{Prove}{$P_i$}}
                \State $RuleProved \gets \text{false}$
                \State \textbf{break}
            \EndIf
        \EndFor
        \If{$RuleProved$}
            \State $WM \gets WM \cup \{G\}$
            \Return \text{true}
        \EndIf
    \EndFor
    \Return \text{false}
\EndFunction
\end{algorithmic}
\end{algorithm}

\section{Подсистемы ЭС}

\subsection{Подсистема объяснения}
Ключевое требование к ЭС --- её прозрачность для пользователя (Explainable AI). Подсистема объяснения обрабатывает два типа запросов:
\begin{itemize}
    \item \textbf{КАК (HOW) получен результат?} Система выдает трассировку сработавших правил (chain of reasoning).
    \item \textbf{ПОЧЕМУ (WHY) система задает вопрос?} Система показывает текущую доказываемую гипотезу (правило, которое она пытается применить).
\end{itemize}

\subsection{Подсистема общения (Интерфейс)}
Предназначена для двустороннего обмена информацией. Исторически интерфейсы ЭС использовали ограниченный естественный язык (NLP) или меню. Основные функции: парсинг ввода пользователя, генерация человекочитаемых вопросов, визуализация данных (например, вероятностей или коэффициентов уверенности).

\subsection{Подсистема приобретения знаний}
Проблема «узкого места» (knowledge acquisition bottleneck) --- самая сложная часть создания ЭС. Подсистема приобретения знаний может быть:
\begin{itemize}
    \item \textit{Пассивной}: Редакторы БЗ с проверкой синтаксиса и семантики (поиск противоречивых или избыточных правил).
    \item \textit{Активной}: Индуктивный вывод правил из примеров (деревья решений, алгоритмы ID3, C4.5), методы машинного обучения.
\end{itemize}

\section{Жизненный цикл экспертной системы}

Жизненный цикл ЭС отличается от классической каскадной модели ПО (waterfall) ввиду высокой неопределенности требований и знаний. Обычно используется \textbf{спиральная модель (быстрое прототипирование)}. Согласно методологии Бьюкенена (Buchanan), выделяют 6 основных этапов:

\begin{enumerate}
    \item \textbf{Идентификация (Identification)} \\
    Определение проблемной области, целей разработки, ресурсов (вычислительных, финансовых), участников (инженер по знаниям, эксперт-доменщик). Проверка того, подходит ли задача для ЭС (необходимость символьных вычислений, отсутствие строгого алгоритмического решения).

    \item \textbf{Концептуализация (Conceptualization)} \\
    Содержательный анализ предметной области. Выделение ключевых понятий, отношений, характеристик данных. Итогом является модель предметной области на неформальном уровне (глоссарий, графы связей).

    \item \textbf{Формализация (Formalization)} \\
    Выбор способа представления знаний (продукции, фреймы, семантические сети, логика предикатов). Выбор инструментальных средств (оболочки ЭС, языки LISP, Prolog, CLIPS).

    \item \textbf{Выполнение / Реализация (Implementation)} \\
    Программирование прототипа. Наполнение БЗ правилами, полученными от эксперта. Создание работающей версии (часто начиная с «микропрототипа» для узкого набора задач).

    \item \textbf{Тестирование (Testing)} \\
    Оценка работы прототипа экспертом (аудит компетенции). Проверка на исторических данных, выявление ошибок в логическом выводе, неполноты или противоречивости базы знаний.

    \item \textbf{Опытная эксплуатация и модификация (Evolution/Maintenance)} \\
    Внедрение ЭС в реальные бизнес-процессы. Система постоянно дообучается, добавляются новые правила. ЭС без регулярного обновления БЗ быстро устаревает (degradation of knowledge).
\end{enumerate}

\end{document}
